%!TEX root = ../main.tex
%=========================================================

\section{Background}
\label{sec:background}
\michal{TODO: mention ENS records (can be very brief)}
%\michal{Need a good introduction on Kademlia here. Distances/buckets etc. It's required to explain the ticket table, routing etc.}
%\etienne{How different is the Ethereum Kademlia from the original published version? I can assume there have been changes in so many years. Is there a reference document for the Ethereum's Kademlia I can read?}
%\sergi{From what i remember main difference are num of buckets and use of log distance.  I don't remember any document but Felix will help.}
%\sergi{We should also add how kademlia is used for discovery (when lookups are performed, lookup buffer use,  pool of connections, random lookups, etc}

%Ethereum is a blockchain platform for distributed computing and was the first major blockchain to support Turing-complete scripting via smart contracts, which are expressed as opcodes specific to the Ethereum Virtual Machine (EVM) [55]. Ethereum also operates as a cryptocurrency by supporting the transfer of tokens called ether between Ethereum accounts. In addition to its use as a fiat currency, ether is also used to incentivize Ethereum nodes to perform dis- tributed computation. Ether can also be converted to gas, which is used as a fee for executing transactions and smart contracts, and effectively mitigates spam. 
%\sergi{Moved to intro.}

\mk{The i-th bucket of a node contains addresses that have a common prefix of length i with the node’s address. Each bucket contains at most k addresses. }
Ethereum uses a Distributed Hash Table (DHT) for P2P communication and implementing higher level protocols. Below, we present an overview of the system and its basic operations. 

\subsection{Ethereum DHT}
\michal{Borrowed from \href{https://arxiv.org/pdf/1908.10141.pdf}{this paper}}
Ethereum DHT is based on Kademlia - a UDP-based, peer-to-peer distributed hash table (DHT) that is used to locate data stord in a decentralized way~\cite{maymounkov2002kademlia}. Each node is uniquely identified by its randomly generated node ID. A node ID in Ethereum is a marshaled version of node's public key. It is easy to generate many different identities but computationally infeasible to generate a specific one. A data item stored in the DHT is also found by its key, which is simply a hash of the data itself. Node IDs and data keys share the same representation; in the following we use the term ID for both.
Ethereum DHT stores data at nodes whose node ID is the “closest” to the data’s key. Closeness is defined as the bitwise XOR between two IDs, taken as an integer value, i.e.,$d(x, y) = \textit{log}_2 \lfloor x \oplus y \rfloor$.


A DHT node stores its known peers in a routing table divided into buckets which partition the known network based on the local node’s ID (\Cref{fig:kademlia}). Every bucket stores up to $k$ neighbors. Bucket $i$ stores nodes whose distance is in $[2i, 2i+1)$, which effectively corresponds to the length of the common prefix between two node IDs. The routing table stores a detailed view of peers close to the node and a less detailed view of parts of the network located further away.

%Items in the DHT and nodes on the network are located by \emph{lookups}. A lookup successively queries nodes that are closest to the desired target ID (key or node ID). Intuitively, a peer closer to the target ID maintains a more detailed view of the target part of the network and provides the initial node with additional peers to query. For instance, the green node in \Cref{fig:kademlia} looking data ID in \emph{bucket 1} will ask known peers in \emph{bucket 1} that will provide additional peers closer to the data ID. Storing data follow the same process.  

\begin{figure}
    \includegraphics[width=0.4\textwidth]{img/kademlia}
    \caption{Ethereum DHT routing table.}
    \label{fig:kademlia}
 \end{figure}
 
In order for a newly instantiated DHT node to find peers, it first adds a hard-coded set of bootstrap node IDs to its routing table. Subsequently, when attempting to locate a target node on the network, the node searches its routing table for the $\alpha$ (typically three) peers that are closest to the target node. The node then sends these $\alpha$ nodes a FIND\_NODE message that specifies the target node ID, and each peer responds with the list of k nodes from its own routing table that are closest to the target. The querying node adds any new node information (i.e., node ID, IP Address, UDP/TCP ports) it discovers through this process to its routing table, and then iteratively repeats the process until it converges on the target node.
 
 Ethereum builds higher level protocols on top of the DHT. RLPx provides secure, point-to-point communication secure by assymetric cryptography. Similarly to TLS, RLPx performs a shared secret derivation authenticated by public/private key pairs of the nodes. The protocol ensures authenticity and integrity of exchanged messages. Furthermore, DEVp2p protocol is used for session management on top of secured channels provided by RLPx. 
 
 The Ethereum DHT is shared across multiple applications.  This includes the Ethereum main net, multiple Ethereum testnets (Rinkeby, Ropsten) or alternative cryptocurrencies (Pirl, Musicoin). In 2018, the Ethereum DHT nodes operated a total of 4,076 different applications~\cite{kim2018measuring}. Currently, a node trying to find application specific peers performs a random walk through the DHT, establishes connections to encountered nodes and verifies the supported protocols at the application level. 

- Existing Discv4 / Eth client properties
- Security challenges including sybil attacks and eclipsing

\subsection{Discovering Peers with Discv4}

In the current discovery system, nodes maintain a discovery table whose sole purpose is to maintain a set of known-peers. The structure of the discovery table is identical to a routing table-each node organises known-nodes according to their distance from itself. A node populates its routing table by attempting to locate a random target node ID and store the nodes encountered during the process in its discovery table, \ie a random walk. When the node locates the closest peers to a random target, the walk is complete. 

Once a node starts populating its discovery table with peers, the node starts establishing an individual application-level session (\ie handshake) with each known-peer in the table in order to identify their subprotocol. In case, the node discovers that a peer is part of a target subprotocol, then the node can send a connection request to the peer (see Section~\ref{sec:connections}) after the handshake; otherwise, the handshake session is terminated. The process of looking for peers is suspended when a node establishes a sufficient number of outbound connections.

This brute-force discovery approach involving random walks and handshakes has several drawbacks. First of all, handshakes with peers that are not running a target subprotocol are wasteful and costs unnecessary overhead. More importantly, random walks can take long time to discover peers that are part of unpopular subprotocols (\ie subprotocols for which there are a small number of peers). 

\subsection{Application-level Connections}
\label{sec:connections}

DEVp2p is responsible for managing the connections to other peers which form the overlay P2P network. For instance,  the Go-ethereum client by default establishes a total of 50  connections to other Ethereum nodes.  Of these 50 slots, around two-thirds (\ie 34) are reserved for inbound connections,  while the remaining 16 are allocated for outbound connections. Inbound connections are initiated by remote peers of a node with a connection request, and the node can decide whether to accept or reject the request. When all the incoming connection slots of a node are full, the node rejects the incoming connection. If a slot is free, a node is free to accept a connection request, although, it is not recommended for simply accepting any request for security reasons as we explain below.

The nodes must carefully select which peers to connect with in order to avoid eclipse attacks. In an eclipse attack, an adversary monopolises all the connections of a victim node with its Sybils, \ie eclipsing the connections, effectively controlling the victim's view of the network. Once it eclipses a victim, an adversary can then feed the victim custom protocol messages as part of a follow-up attack such as double spending and stubborn mining~\cite{henningsen2019eclipsing}.

\ramin{Below and above pargraph have a lot of overlap. Merge?}

Avoiding eclipsing requires careful selection of peers to fill the 50 slots. The set of connections require protection from Eclipsing attacks where an adversary uses a set of Sybil nodes to monopolize the connections of a victim, isolating from the rest of the network. Inbound connections are chosen from incoming connection requests from peers. Enforcing IP diversity in the chosen inbound peers is recommended by existing studies~\cite{henningsen2019eclipsing, marcus2018low}. Outbound connections are chosen locally and picking a random subset from a large list of known-nodes is shown to be effective against eclipsing attacks. 


%In an eclipse attack, an
%adversary monopolizes the connections of a victim, effectively
%filtering the victim’s view of the blockchain.  Eclipse attacks
%enable a variety of follow-up attacks such as double spending
%and stubborn mining~\cite{}.
%\michal{ditto}
